\chapter{Reproducibility Guide}\label{app:reproducibility}

\section{Where everything is}
\begin{itemize}
\item \textbf{Code, configurations and raw logs:} the project repository \url{https://github.com/Unavailable4u/meeting-loss} (commit \code{\repocommit} for the analysis in this thesis). Raw run logs are in \code{results/raw/} (one JSON-lines file per run, with a header that records the configuration, its hash, the git commit and whether the working tree was clean); processed tables are in \code{results/processed/}.
\item \textbf{Locks:} tags \code{prereg-v1} (\code{3378d36}), \code{prereg-v2} (\code{cc2830c}) and \code{prereg-v3} (\code{a49e9e3}). No source, configuration or hypothesis file locked by a tag changed between the tag and the final commit, except where the drift log records an additive change.
\item \textbf{This thesis:} the \code{scripts/} folder contains the Python code that produces every figure (\code{fig\_study1.py}, \code{fig\_study23.py}) and every table (\code{make\_tables.py}) directly from the raw logs, using the repository's own scoring and bootstrap functions (\code{src/sweep.py}, \code{src/prereg3.py}). The \LaTeX{} tables in \code{tables/} and the PDF figures in \code{figures/} are the outputs of these scripts, so the numbers in the text can be re-checked.
\end{itemize}

\section{Regenerating figures and tables}
\begin{lstlisting}
git clone https://github.com/Unavailable4u/meeting-loss.git
export MEETING_LOSS_REPO=$PWD/meeting-loss
pip install numpy pandas pyyaml matplotlib
cd thesis/scripts
python make_tables.py       # writes ../tables/*.tex
python fig_study1.py        # writes ../figures/fig_s1_*.pdf, fig_tuning_noise, fig_dose, ...
python fig_study23.py       # writes ../figures/fig_s2_*.pdf, fig_s3_*.pdf, fig_metrics.pdf
cd .. && latexmk -pdf main.tex
\end{lstlisting}
The figure scripts import the repository's \code{src} package; \code{src/prereg3.py} needs only NumPy and the standard library at analysis time. Re-running the experiments themselves requires PyTorch (CPU is sufficient) and the commands in \code{scripts/README.md} of the repository; the 260 confirmatory runs of Study~3 took 15\,856 seconds ($4.4$ hours) on four worker processes, as logged.

\section{Seed ranges}
\begin{table}[h]
\centering\small
\caption{Seed ranges by purpose. No seed is used for both tuning and evaluation, and the exploratory seeds are never reused in a confirmatory test.}
\begin{tabular}{@{}l l@{}}
\toprule
Purpose & Seeds\\
\midrule
Tuning (selection), tags \code{t1}, \code{t2}, \code{t5}--\code{t8} & 1000/1001, 2000/2001, 5000/5001, \dots, 8000/8001\\
Tuning of B48, S30T, SBT (\code{t9}, \code{t10}, \code{t11}) & 9000/9001, 9100/9101, 9200/9201\\
Study 1 (confirmatory) & 0--19 (A, D, B, A2); 0--9 (other groups)\\
Concentration, geometry, pilot (exploratory) & 500--502; 600--602\\
Dose-response blocks & 1100--1109, 1110--1119, 1120--1129\\
Timing and small-batch looks & 1200--1204\\
Study 2 (confirmatory) & 1300--1319\\
Headroom and look 2 (exploratory) & 2100--2109, 2110--2119\\
Study 3 (confirmatory) & 2300--2319\\
\bottomrule
\end{tabular}
\end{table}

\section{Status of verification at the time of writing}
\begin{itemize}
\item \textbf{Numbers.} Every number in the thesis tables and figures was recomputed from the raw logs by the scripts above; the recomputed values reproduce the values in the repository's experiment log for Study 1 (E-008), Study 2 (E-022) and Study 3. At the time of writing, the repository's own documents (\code{docs/HYPOTHESES.md}, \code{docs/FINDINGS.md}, \code{docs/EXPERIMENT\_LOG.md}) had not yet recorded the verdicts of H8--H10 (the status lines still read ``open''); the processed result file \code{results/processed/prereg\_v3.csv} and the analysis script \code{scripts/analyze\_prereg3.py} contain them, and the recomputation here gives the same verdicts. The repository documents should be updated before they are cited.
\item \textbf{Unit tests.} The repository contains 24 test files with 223 test functions (counted, not run, in the environment used to write this thesis, which had no PyTorch). Running them is a necessary step before submission: \code{pytest tests/}.
\item \textbf{Bibliography.} The entries for Jiang et al.\ (Selective-Backprop), Kaddour et al., Shibuya et al., Lee et al.\ (DTP), Kawaguchi and Lu, Bartunov et al.\ and Meulemans et al.\ were checked against web sources while writing this thesis. The remaining entries were written from memory, without page numbers, and must be checked against the publisher pages. The literature review is a focused one and the systematic related-work check listed in \code{docs/LITERATURE.md} remains open (\cref{sec:future}).
\end{itemize}
